\documentclass[8pt, a4paper]{extarticle}
\usepackage[
    a4paper, 
    landscape, 
    margin=0.2in,
    footskip=0.1in,
    headsep=0in,
    headheight=0in
]{geometry}
\usepackage{graphicx}
\usepackage{ctex}
\usepackage{multicol}
\usepackage{amsmath,amssymb,amsfonts,amsthm}
\usepackage{enumitem}
\usepackage{pifont}

\usepackage{tabularx}
\usepackage{booktabs}
\usepackage{multirow}
\renewcommand{\tabularxcolumn}[1]{m{#1}}
\newcolumntype{C}[1]{>{\hsize=#1\hsize\centering\arraybackslash}X}

\usepackage{tikz}
\usetikzlibrary{arrows.meta, positioning, shapes, fit, backgrounds}

\usepackage{fix-cm}
\DeclareMathSizes{7}{7}{4}{4}

\setlength{\columnsep}{0.5cm}
\setlength{\parindent}{0pt}
\setlist{noitemsep}
\setlength{\jot}{0.5pt}
\setcounter{secnumdepth}{0}
\pagestyle{empty}

\usepackage[tiny,compact]{titlesec}
\titleformat{\section}{\bfseries\normalsize}{\thesection}{0.5em}{}
\titlespacing*{\section}{0pt}{1ex}{0.5ex}
\titleformat{\subsection}{\bfseries\small}{\thesubsection}{0.4em}{}
\titlespacing*{\subsection}{0pt}{0.5ex}{0.2ex}
\titleformat{\subsubsection}{\normalfont\small\bfseries\itshape}{\thesubsubsection}{0.3em}{}
\titlespacing*{\subsubsection}{0pt}{0.4ex plus .1ex minus .1ex}{0.2ex}
\titleformat{\paragraph}[runin]{\normalfont\small\bfseries}{\theparagraph}{0.5em}{}[.]
\titlespacing*{\paragraph}
  {0pt}
  {0.3ex plus 0.1ex minus 0.1ex}
  {0.5em} 

\newcommand*{\dif}{\mathop{}\!\mathrm{d}}
\DeclareMathOperator{\sinc}{sinc}
\DeclareMathOperator{\Sa}{Sa}
\DeclareMathOperator{\rect}{rect}
\DeclareMathOperator{\pv}{p.v.}
\DeclareMathOperator{\sgn}{sgn}
\DeclareMathOperator{\E}{E}
\DeclareMathOperator{\prob}{P}
\DeclareMathOperator{\fourier}{\mathcal{F}}
\DeclareMathOperator{\bit}{bit}
\DeclareMathOperator{\erf}{erf}
\DeclareMathOperator{\erfc}{erfc}

\makeatletter
\renewcommand{\maketitle}{
  \begin{center}
    \vspace{-2.5em} 
    
    {\Large \bfseries \@title} \\[0.4ex] 
    
    {\scriptsize \@author \quad \text{|} \quad \@date} 
    
    \vspace{0.2ex}
    \hrule 
    
    \vspace{-1.0em} 
  \end{center}
}
\makeatother

\title{奇襲通原\ding{43}}
\author{luisleee}
\date{\today}

\begin{document}
\small
\begin{multicols}{5}

\maketitle
\section{基本约定}
除非特别说明，下面的叙述全部关于复信号。记号上虚数单位用\(j\)表示，用\({(\cdot)}^*\)表示复共轭，上划线\(\overline{(\cdot)}\)表示时间平均。复内积都保持第一个位置线性，第二个位置共轭线性。卷积用\(\ast\)表示。傅里叶变换采用模拟频率\(f\)的形式，用\(\fourier[\cdot]\)表示，变换对用\(\Leftrightarrow\)表示。柯西主值用\(\pv\)表示。对随机变量的积分和极限都是均方意义下的。\(\log\)表示以\(2\)为底的对数。
\subsection{常用函数性质}
单位冲激信号\((f\ast \delta)(t)=f(t)\)

阶跃函数\(u(t) = 1/2(1+\sgn(t))\)

采样函数\(\Sa(t) =\sin t/t,\, \Sa(0)=1\)

\(\sinc\)脉冲\(\sinc(t) = \Sa(\pi t)\)

单位矩形脉冲\(\rect(t) = u(t-1/2)-u(t+1/2)\)

\(n\)阶Bessel函数
\[
J_n(\beta)=\frac{1}{2\pi}\int_{-\pi}^{\pi}e^{j(\beta\sin\theta-n\theta)}\dif \theta
\]

误差函数
\[
\erf(x)=\frac{2}{\sqrt{\pi}}\int_{0}^{x}e^{-t^2}\dif t
\]

互补误差函数\(\erfc(x)=1-\erf(x)\)
\section{通信系统模型}
\textbf{信源} $\to$ [信源编码 $\to$ 信道编码] $\to$ \textbf{信道} $\to$ [信道译码 $\to$ 信源译码] $\to$ \textbf{信宿}
\section{确定信号分析}
\subsection{基本概念}
信号\(x(t)\)的时间平均
\[
\overline{x(t)} = \lim_{T\to\infty}\frac{1}{T}\int_{-T/2}^{+T/2}x(t)\dif t
\]
能量
\[
E_x = \int_{-\infty}^{+\infty}|x(t)|^2\dif t
\]
功率
\[
P_x = \overline{|x(t)|^2}
\]
能量信号\(x(t),y(t)\)的内积（互能量）
\[
E_{xy} =\int_{-\infty}^{+\infty}x(t)y^*(t)\dif t
\]
功率信号\(x(t),y(t)\)的内积（互功率）
\[
P_{xy} = \overline{x(t)y^*(t)}
\]
Cauchy-Schwarz 不等式
\[
\langle x,y\rangle \leq \sqrt{\|x\|^2\|y\|^2}
\]
\subsection{频谱}
周期为\(T\)的信号\(x(t)\)的傅里叶级数
\[
x(t) = \sum_{n=-\infty}^{+\infty}x_n e^{j2\pi nft}, \quad f=\frac{1}{T}
\]
\[
x_n = \frac{1}{T}\int_{-T/2}^{+T/2}x(t)e^{-j2\pi nft}\dif t
\]
能量信号\(x(t)\)的傅里叶变换
\[
X(f)=\int_{-\infty}^{+\infty}x(t)e^{-j2\pi ft}\dif t
\]
\(X(f)\)的逆傅里叶变换
\[
x(t)=\int_{-\infty}^{+\infty}X(f)e^{j2\pi ft}\dif f
\]
能量信号\(x(t)\)的能量谱密度
\[
E_x(f) = |X(f)|^2
\]
能量信号\(x(t),y(t)\)的互能量谱密度
\[
E_{xy}(f) = X(f)Y^*(f)
\]
功率信号\(x(t)\)的功率谱密度，截短信号
\newline
\(x_T(t)=\rect(t/T) x(t)\)
\[
P_x(f) = \lim_{T\to\infty}\frac{1}{T}|X_T(f)|^2
\]
功率信号\(x(t),y(t)\)的互功率谱密度
\[
P_{xy}(f)=\lim_{T\to\infty}\frac{1}{T}X_T(f)Y_T^*(f)
\]

\subsubsection{常用变换对及性质}
面积和原点值
\begin{gather*}
x(0) = \int_{-\infty}^{+\infty}X(f)\dif f\\
X(0) = \int_{-\infty}^{+\infty}x(t)\dif t
\end{gather*}
线性
\(
ax(t)+y(t)\Leftrightarrow aX(f)+Y(f)
\)

共轭
\(
x^*(t) \Leftrightarrow X^*(-f)
\)

奇偶性
\(
x(-t) \Leftrightarrow X(-f)
\)

对偶
\(
X(t) \Leftrightarrow x(-f)
\)

尺度变换
\(
x(at) \Leftrightarrow 1/|a| X\left(f/a\right)
\)

\(\rect\)和\(\sinc\)
\begin{gather*}
\rect(t)\Leftrightarrow \sinc(f)\\
\sinc(t)\Leftrightarrow \rect(f)
\end{gather*}
直流和冲激\(1\Leftrightarrow \delta(f),\quad \delta(t)\Leftrightarrow 1\)

时移和频移
\begin{gather*}
x(t-t_0) \Leftrightarrow X(f)e^{-j2\pi ft_0}\\
x(t)e^{j2\pi f_0 t}\Leftrightarrow X(f-f_0)
\end{gather*}
调制定理
\[
x(t)\cos(2\pi f_0 t) \Leftrightarrow \frac{1}{2}[X(f-f_0)+X(f+f_0)]
\]
微分
\begin{gather*}
\frac{\dif}{\dif t}x(t) \Leftrightarrow j2\pi f \cdot X(f)\\
-j2\pi t \cdot x(t) \Leftrightarrow \frac{\dif}{\dif f}X(f)
\end{gather*}
周期冲激序列
\[
\sum_{n=-\infty}^{+\infty}\delta(t-nT) \Leftrightarrow \frac{1}{T}\sum_{n=-\infty}^{+\infty}\delta \left(f-n/T\right)
\]
卷积定理
\[
\begin{gathered}
(x\ast y)(t) \Leftrightarrow X(f)Y(f)\\
x(t)y(t) \Leftrightarrow (X\ast Y)(f)    
\end{gathered}
\]

内积
\[
\int_{-\infty}^{+\infty}x(t)y^*(t)\dif t = \int_{-\infty}^{+\infty}X(f)Y^*(f)\dif f 
\]
Parseval定理
\[
\int_{-\infty}^{+\infty}|x(t)|^2\dif t = \int_{-\infty}^{+\infty}|X(f)|^2\dif f 
\]
\subsubsection{相关函数和谱}
能量信号\(x(t)\)的自相关函数
\[
R_x(\tau) = \int_{-\infty}^{+\infty}x(t+\tau)x^*(t)\dif t
\]
能量信号\(x(t),y(t)\)的互相关函数
\[
R_{xy}(\tau) = \int_{-\infty}^{+\infty}x(t+\tau)y^*(t)\dif t
\]
功率信号\(x(t)\)的自相关函数
\[
R_x(\tau) = \overline{x(t+\tau)x^*(t)}
\]
功率信号\(x(t),y(t)\)的互相关函数
\[
R_{xy}(\tau) = \overline{x(t+\tau)y^*(t)}
\]
互相关函数的共轭对称性
\[
R_{yx}(\tau)=R_{xy}^*(-\tau)
\]
自相关函数最大值
\[
|R_x(\tau)|\leq R_x(0)
\]
相关定理
\[
R_{xy}(\tau) \Leftrightarrow E_{xy}(f)
\]
Wiener-Khinchin定理
\[
R_{xy}(\tau) \Leftrightarrow P_{xy}(f)
\]
能量守恒
\[
E_x = \int_{-\infty}^{\infty} E_x(f) \dif f = R_x(0)
\]
功率守恒
\[
P_x = \int_{-\infty}^{\infty} P_x(f) \dif f = R_x(0)
\]
\subsection{线性滤波器}
滤波器的冲激响应为\(h(t)\)，\(g(t)=h^*(-t)\)，输入\(x(t)\)，输出
\[
y(t) = (x\ast h)(t)= R_{xg}(t)
\]
频域特性
\[
Y(f)=H(f)X(f)
\]
无失真条件
\[
H(f)=c\cdot e^{-j2\pi ft_0}
\]
通过滤波器后的相关函数
\begin{gather*}
R_y(\tau) = (R_X\ast h \ast g)(\tau)\\
R_{xy}(\tau) = (R_X\ast g)(\tau)\\
R_{yx}(\tau) = (R_X\ast f)(\tau)
\end{gather*}
\subsection{带通信号表示}
实信号\(x(t)\)的希尔伯特变换
\[
\hat{x}(t) = \frac{1}{\pi}\pv\int_{-\infty}^{+\infty}\frac{x(\tau)}{t-\tau}\dif \tau
\]
逆希尔伯特变换
\[
x(t) = -\frac{1}{\pi}\pv\int_{-\infty}^{+\infty}\frac{\hat{x}(\tau)}{t-\tau}\dif \tau
\]
希尔伯特变换的传递函数
\[
H(f) = -j\sgn(f),\quad |H(f)|^2 = 1
\]
正交关系
\[
R_{\hat{x}x}(0) = \int_{-\infty}^{+\infty}\hat{x}(t)x(t)\dif t = 0
\]
实信号\(x(t)\)的解析信号
\[
z(t) =x(t)+j\hat{x}(t)
\]
解析信号的频谱，\(X_+(f)=X(f)u(f)\)表示正频率部分，\(X_-(f)=X(f)u(-f)\)表示正频率部分
\begin{gather*}
Z(f)=2X_+(f)\\
Z^*(-f)=2X_-(f)
\end{gather*}
相关函数关系
\begin{gather*}
R_z(\tau)=2[R_x(\tau)+j\hat{R}_x(\tau)]\\
R_{zz^*}(\tau)=0
\end{gather*}
实信号\(x(t)\)关于参考载波\(\cos(2\pi f_c t +\theta_c)\)的复包络；\(x_c(t)\)为同相分量（I路分量），\(x_s(t)\)为正交分量（Q路分量）
\[
x_L(t) = z(t)e^{-j(2\pi f_c t+\theta_c)} = x_c(t) + jx_s(t)
\]
逆关系
\[
z(t) = x_L(t)e^{j(2\pi f_c t+\theta_c)}
\]
自相关函数
\[
R_{x_L} = 2[R_x(\tau)+j\hat{R}_x(\tau)]e^{-j2\pi f_c \tau}
\]
I/Q调制关系
\begin{align*}
x(t)&=x_c(t)\cos(2\pi f_c t+\theta_c)\\
&-x_s(t)\sin(2\pi f_c t+\theta_c)
\end{align*}
I/Q解调，相位不同步发生旋转
\begin{gather*}
x(t)\cdot 2\cos(2\pi f_c t+\theta) = x_c(t)\cos(\theta_c-\theta) + \cdots\\
x(t)\cdot -2\sin(2\pi f_c t+\theta) = x_s(t)\cos(\theta_c-\theta) + \cdots\\
y_L(t) = x_L(t)e^{j(\theta_c-\theta)}
\end{gather*}
频谱关系
\begin{gather*}
X_L(f) = 2X_+(f+f_c)\\
X_L^*(-f) = 2X_-(f+f_c)\\
X(f) = \frac{1}{2}\left(X_L(f-f_c) + X_L^*(-f-f_c)\right)\\
X_c(f) = X_+(f+f_c) + X_-(f-f_c)\\
X_s(f) = X_+(f+f_c) - X_-(f-f_c)
\end{gather*}
系统等效基带表示
\[
H_e(f) = \frac{1}{2}H_L(f) = H_+(f+f_c)
\]
\section{随机过程}
\subsection{基本概念}
随机过程\(X(t)\)的均值（统计平均），平稳过程与\(t\)无关
\[
m_X(t) = \E[X(t)]
\]
随机过程\(X(t)\)的时间平均，为一个随机变量
\[
\overline{X(t)} = \lim_{T\to \infty}\frac{1}{T}\int_{-T/2}^{+T/2}X(t)\dif t
\]
统计平均和时间平均可交换（在可积的条件下）
\[
\E[\overline{X(t)}] = \overline{\E[X(t)]}
\]
随机过程\(X(t)\)的平均功率
\[
P_X=\overline{\E[|X(t)|^2]}
\]
随机过程\(X(t)\)的自相关函数
\[
R_X(t_1,t_2) = \E[X(t_1)X^*(t_2)]
\]
随机过程\(X(t),Y(t)\)的互相关函数
\[
R_{XY}(t_1,t_2) = \E[X(t_1)Y^*(t_2)]
\]
随机过程\(X(t)\)的平均自相关函数
\[
\overline{R}_X(\tau) = \overline{R_X(t+\tau,t)}
\]
随机过程\(X(t),Y(t)\)的平均互相关函数
\[
\overline{R}_{XY}(\tau) = \overline{R_{XY}(t+\tau,t)}
\]
平稳过程\(X(t)\)的自相关函数
\[
R_{X}(\tau) = \E[X(t+\tau)X^*(t)]
\]
平稳过程\(X(t),Y(t)\)的互相关函数
\[
R_{XY}(\tau) = \E[X(t+\tau)Y^*(t)]
\]
平稳过程\(X(t)\)是遍历的即
\begin{gather*}
\prob[X(t)=m_X]=1\\
\prob[\overline{X(t+\tau)X^*(t)}=R_X(\tau)]=1
\end{gather*}
平稳过程\(X(t)\)的功率谱密度，截短过程\(X_T(t)=\rect(t/T) X(t)\)
\[
P_X(f) = \lim_{T\to \infty}\frac{1}{T}\E\{|\fourier[X_T(t)]|^2\}
\]
平稳过程\(X(t),Y(t)\)的互功率谱密度
\[
P_{XY}(f) = \lim_{T\to \infty}\frac{1}{T}\E\{\fourier[X_T(t)]{\fourier[Y_T(t)]}^*\}
\]
Wiener-Khinchin定理
\[
\overline{R}_{XY}(\tau) \Leftrightarrow P_{XY}(f)
\]
随机过程\(X(t)\)通过滤波器，输出\(Y(t)\)
\begin{gather*}
P_Y(f) = |H(f)|^2P_X(f)\\
P_{XY}(f) = H^*(f)P_X(f)\\
P_{YX}(f) = H(f)P_X(f)
\end{gather*}
\subsection{AWGN}
功率谱密度
\(
P_{n_w}(f) = N_0/2
\)

自相关函数
\(
R_{n_w}(\tau) =N_0/2\delta(\tau)
\)

通过滤波器
\begin{gather*}
P_n(f)=\frac{N_0}{2}|H(f)|^2\\
P_n = \frac{N_0}{2}E_h
\end{gather*}

\(n_w(t)\)和确定信号\(x(t)\)的内积

\(Z_x\sim N(0,N_0E_x/2)\)，\(\E[Z_x Z_y] = N_0E_{xy}/2\)

AWGN信道模型；发送信号\(s(t)\)，接收信号\(r(t)\)，噪声\(n_w(t)\)独立
\[
r(t)=s(t)+n_w(t)
\]
\subsection{匹配滤波}
滤波器在 \( t = t_0 \) 时刻的输出信号
\[
y(t_0) = (s\ast h)(t_0) + (n_w\ast h)(t_0) = s_0 + Z
\]
\(Z\)的平均功率
\(
P_n=N_0/2 E_h
\)

采样时刻信噪比
\(
\gamma = |s_0|^2/P_n
\)

最大值条件
\[
|s_0|^2 = \langle h(t),s^*(t_0-t) \rangle \leq E_h E_s
\]
最大信噪比
\(
\gamma_{\max} = 2E_s/N_0
\)

对信号\(s(t)\)匹配的匹配滤波器
\[
\begin{gathered}
h(t)=Ks^*(t_0-t)\\
H(f)=KS^*(f)e^{-j2\pi ft_0}
\end{gathered}
\]

\section{模拟通信系统}
\subsection{DSB-SC}
\subsubsection{调制}
已调信号
\begin{gather*}
s(t)=A_c m(t)\cos(2\pi f_c t)\\
S(f)=\frac{A_c}{2}[M(f-f_c)+M(f+f_c)]\\
P_s(f)=\frac{A_c^2}{4}[P_m(f-f_c)+P_m(f+f_c)]
\end{gather*}
\subsubsection{相干解调}
\(\phi_c\)为发端载波初相，\(\phi\)为收端相位
\[
y_o(t)=A_c m(t)\cos(\phi_c -\phi)
\]
\subsubsection{载波同步}
\paragraph{插入导频}
等效为\(m(t)\)叠加直流\(A_p/A_c\)，提取时用窄带滤波或者锁相环
\begin{align*}
s(t)&=A_c m(t)\cos(2\pi f_c t+\phi_c) \\
&+ A_p\cos(2\pi f_c t+\phi_c)
\end{align*}
锁相环用\(\tilde{Q}(t)\)负反馈，提取载波，锁定时输出
\[
\tilde{I}(t)\approx A_c m(t)+A_p
\]
\paragraph{平方环}
平方得
\[
A_c^2m^2(t)\cos^2[2\pi f_c t +\phi_c]
\]
锁相环提取\(\cos(4\pi f_c t +2\phi_c)\)后二分频得载波
\[
\pm\cos(2\pi f_c t +\phi_c)
\]
\paragraph{科斯塔斯环}
用\(\tilde{I}(t)\tilde{Q}(t)\)负反馈，锁定时解调输出
\[
\tilde{I}(t)\approx \pm A_c m(t)
\]
\subsection{AM}
\subsubsection{调制}
调幅系数\(a\leq 1\)，最大幅度归一化的\(m_n(t)=m(t)/|m(t)|_{\max}\)。两种方法为直流偏置后DSB调制、DSB-SC叠加载波。
\begin{align*}
s(t) &=A_c[1+am_n(t)]\cos(2\pi f_c t)\\
&=[A_c+A'm(t)]\cos(2\pi f_c t)\\
&=A_c\cos(2\pi f_c t)+A'm(t)\cos(2\pi f_c t)
\end{align*}
频谱特性
\begin{align*}
S(f)&=\frac{A'}{2}[M(f-f_c)+M(f+f_c)]\\
&+\frac{A_c}{2}[\delta(f-f_c)+\delta(f+f_c)]\\
P_s(f)&=\frac{{(A')}^2}{4}[P_m(f-f_c)+P_m(f+f_c)]\\
&+\frac{A_c^2}{4}[\delta(f-f_c)+\delta(f+f_c)]
\end{align*}
\subsubsection{解调}
非相干解调，包络检波得\(A_c+A'm(t)\)，隔直流得\(A'm(t)\)。
\subsubsection{调制效率}
调制效率，\(P_{m_n}\)为\(m_n(t)\)的功率
\[
\eta = \frac{{(A')^2}P_m}{A_c^2+{(A')}^2P_m} =\frac{a^2P_{m_n}}{1+a^2P_{m_n}}
\]
峰均比（PAPR）
\[
C_m=\frac{|m(t)|^2_{\max}}{\overline{m^2(t)}}=P_{m_n}^{-1}
\]
调制效率、峰均比、调幅系数关系
\[
\eta = \frac{1}{1+C_m/a^2}
\]
\subsection{SSB}
上下边带SSB信号
\begin{align*}
s_{USB}(t)&=\frac{A_c}{2}[m(t)\cos(2\pi f_c t)\\
&-\hat{m}(t)\sin(2\pi f_c t)]\\
s_{LSB}(t)&=\frac{A_c}{2}[m(t)\cos(2\pi f_c t)\\
&+\hat{m}(t)\sin(2\pi f_c t)]
\end{align*}
\subsection{VSB}
带通滤波的等效基带传递函数\(H_e(f)\)满足条件
\[
H_e(f)+H_e^*(-f)=1
\]
带通滤波\(H(f)\)满足条件
\[
H(f+f_c)+H(f-f_c)=1,\quad |f|\leq f_c
\]
\subsection{PM/FM}
PM/FM已调信号；\(K_p\)调相灵敏度，\(K_f\)调频灵敏度
\begin{gather*}
s_{PM} =A_c\cos\left[2\pi f_c t + K_p m(t)\right]\\
s_{FM} =A_c\cos\left[2\pi f_c t + 2\pi K_f \int_{-\infty}^{t}m(\tau)\dif \tau\right]
\end{gather*}
瞬时频偏关系
\begin{gather*}
\theta_p(t) = K_p m(t)\\
\frac{1}{2\pi}\theta_m'(t)=K_f m(t) 
\end{gather*}
FM最大频偏
\[
\Delta f_{\max} =K_f|m(t)|_{\max}
\]
调频指数，\(m(t)\)最高频率\(W\)
\[
\beta_f = \Delta f_{\max}/W
\]
最大幅度归一化信号\(m_n(t)\)的形式，\(\tilde{m}(t)=2\pi W\int_{-\infty}^{t}m_n(\tau)\dif \tau\)
\[
s_{FM}=A_c\cos\left[2\pi f_c t + \beta_f \tilde{m}(t)\right]
\]
单音信号\(m_n(t)=\cos(2\pi f_m t)\)的调频，此时\(W=f_m\)
\[
s_{FM}(t)=A_c\cos\left[2\pi f_c t + 2\pi \beta_f \sin(2\pi f_m t)\right]
\]
单音调频的复包络
\begin{align*}
s_L(t)&=A_c e^{j\beta_f\sin(2\pi f_m t)}\\
&=A_c\sum_{n=-\infty}^{+\infty}J_n(\beta_f)e^{j2\pi n f_m t}
\end{align*}
单音调频复包络功率谱密度
\[
P_{s_L}(f)=A_c\sum_{n=-\infty}^{+\infty}J_n^2(\beta_f)\delta(f-nf_m)
\]
带宽Carson公式
\[
B\approx 2(\beta_f+1)W=2(\Delta f_{\max} +W)
\]
窄带调频，\(|\theta(t)=\beta_f m_n(t)|\)非常小
\begin{align*}
s_L(t) &=A_c e^{j\beta_f \tilde{m}(t)}\\
&\approx A_c[1+j\beta_f \tilde{m}(t)]\\
s(t) &\approx A_c\cos(2\pi f_c t)-A_c\beta_f \tilde{m}(t)\sin(2\pi f_c t)
\end{align*}
\subsubsection{调制}
直接调频采用VCO。间接调频先窄带调频，然后\(n\)倍频增大调频指数\(f'=nf,\beta_f'=n\beta_f\)。
\subsubsection{解调}
\paragraph{微分包络检波法}\(v(t)=\frac{1}{2\pi}\frac{\dif}{\dif t}s(t)\)
\begin{align*}
v(t) &=-[f_c+K_f m(t)]\\
&\cdot \sin\left[2\pi f_c t+2\pi K_f \int_{-\infty}^{t}m(\tau)\dif \tau\right]
\end{align*}
\(v(t)\)包络检波后得\(f_c+K_f m(t)\)，隔直流为\(K_f m(t)\)
\paragraph{锁相鉴频法}锁相环锁定时，VCO调频灵敏度相同
\[
2\pi K_f\int_{-\infty}^{t}m(\tau)\dif \tau \approx 2\pi K_f\int_{-\infty}^{t}v(\tau)\dif \tau
\]
\subsection{抗噪声性能对比}
输入信噪比
\(
SNR_i = P_R/N_0 B
\)

相干解调相位不理想同步
\[
SNR_o' = SNR_o \cos^2\theta_e
\]
\begin{tabularx}{\linewidth}{|C{0.8}|C{0.8}|C{1.0}|C{1.4}|}
    \toprule
    调制 & 带宽$B$ & 解调 & \(SNR_o\) \\
    \midrule
    DSB-SC & $2W$  & 相干解调 & $\frac{P_R}{N_0 W}$ \\
    \hline
    \multirow{2}{*}{AM} & \multirow{2}{*}{$2W$}  & 相干解调 & $\eta\frac{P_R}{N_0 W}$ \\
    \cline{3-4}
    &   & 包络检波 & ${\eta\frac{P_R}{N_0 W}}^{*}$ \\
    \hline
    SSB & $W$  & 相干解调 & $\frac{P_R}{N_0 W}$ \\
    \hline
    FM & $2(\beta_f+1)W$  & 鉴频器 & ${\frac{3\beta_f^2}{C_m}\frac{P_R}{N_0 W}}^{*}$ \\
    \bottomrule
    \multicolumn{4}{@{}l@{}}{\tiny * 大信噪比下近似。} 
\end{tabularx}
\subsection{频分复用}
调制一般先采用SSB,DSB，频带互不交叠，再FM调制后传输。解复用使用带通滤波。
\section{数字基带传输}
\subsection{基本速率}
比特速率\(R_b\)和二进制符号间隔\(R_b\)关系
\[
R_b T_b = 1 \bit
\]
比特速率\(R_b\)和符号速率\(R_s\)关系，\(M\)进制符号
\[
R_s=\frac{R_b}{\log M}
\]
比特错误率\(P_b\)与符号错误率\(P_s\)关系
\[
\frac{R_s}{\log M}\leq P_b \leq P_s
\]
\subsection{PAM}
\(M\)进制幅度序列，幅度\(a_n\)有\(M\)种取值，\(M\)进制符号间隔\(T_s\)
\[
\sum_{n=-\infty}^{+\infty}a_n \delta(t-nT_s)
\]
PAM信号一般表达式，发送滤波器冲激响应\(g_T(t)\)
\[
s(t)=\sum_{n=-\infty}^{+\infty}a_n g_T(t-nT_s)
\]
功率谱
\[
P_s(f)=\frac{1}{T_s}\sum_{n=-\infty}^{+\infty}R_a(n)e^{j2\pi nfT_s}|G_T(f)|^2
\]
\(\{a_n\}\)是不相关序列，均值\(m_a\)，方差\(\sigma_a^2\)
\[
\begin{aligned}
P_s(f)&=\frac{\sigma_a^2}{T}|G_T(f)|^2\\
&+\frac{m_a^2}{T_s^2}\sum_{k}\left|G_T\left(\frac{k}{T_s}\right)\right|^2\delta\left(f-\frac{k}{T_s}\right)
\end{aligned}
\]
\(\{a_n\}\)为零均值不相关序列时无离散谱
\subsection{常用码型}
单极性码，可能的幅值为\(A,0\)；双极性码，可能的幅值为\(A,-A\)。差分码\(b_n=a_n\oplus b_{n-1}\)，原序列独立等概率则仍然独立等概率。

矩形不归零脉冲，占空比\(100\%\)
\[
\begin{gathered}
g_T(t)=A_b(u(t-T_b)-u(t))\\
G_T(f)=A_b T_b\sinc(fT_b) e^{-j\pi fT_b}
\end{gathered}
\]
矩形归零脉冲，占空比\(50\%\)
\[
\begin{gathered}
g_T(t)=u(t-T_b/2)-u(t)\\
G_T(f)=A_b T_b/2\sinc(fT_b/2) e^{-j\pi fT_b/2}
\end{gathered}
\]
\subsection{常用线路码型}
AMI码：传号“1”交替反转。
\[
R_m(0)=1/2,\quad R_m(\pm 1) = -1/4
\]
功率谱（占空比\(50\%\)）
\[
P(f)=\frac{T_b A_b^2}{4}\sin^2(\pi fT_b)\sinc^2(fT_b/2)
\]

HDB\textsubscript{3}码：“1”交替反转，将四连“0”替换为取代节，插入破坏符号\(V\)。

CMI码：将“0”编码为“01”，将“1”交替编码为“11”、“00”。

数字双相码：将“0”编码为“01”，将“1”编码为“10”。
\subsection{噪声信道下的接收}
双极性信号
\(
s(t)=\pm g(t)
\)

匹配滤波器
\(
h(t)=g(t_0-t)
\)

平均比特能量
\(
E_b=E_g
\)

判决门限\(V_T=0\)

平均误比特率，只与比特信噪比\(E_b/N_0\)有关
\[
P_b=\frac{1}{2}\erfc\left(\sqrt{\frac{E_b}{N_0}}\right)
\]

单极性信号\(s(t)=g(t)\text{或}0\)

平均比特能量\(E_b=E_g/2\)

判决门限\(V_T=E_g/2\)

平均误比特率
\[
P_b=\frac{1}{2}\erfc\left(\sqrt{\frac{E_b}{2N_0}}\right)
\]
\subsection{PAM有限带信道传输}
PAM基带传输系统模型 \(g_T(t)\to c(t) \to +n_w(t) \to g_R(t)\)

输出信号
\[
\begin{gathered}
y(t)=\sum_{n=-\infty}^{+\infty}a_n x(t-nT_s) + \gamma(t)\\
x(t)=(g_T\ast c\ast g_R)(t)\\
\gamma(t)=(n_w\ast g_R)(t)
\end{gathered}
\]
瞬时采样值
\[
y_m = x_0a_m + \sum_{n\neq m}a_n x_{m-n} + \gamma_m
\]
无ISI传输的Nyquist准则，假设理想基带信道
\[
\begin{gathered}
x(nT_s)=\delta(n)\\
\Leftrightarrow \sum_{n=-\infty}^{+\infty}X(f+\frac{n}{T_s})=T_s
\end{gathered}
\]

\end{multicols}
\end{document}
